\documentclass[11pt,a4paper]{article}
\usepackage{amsmath,amssymb,amsthm,mathtools,booktabs,microtype}
\usepackage[margin=25mm]{geometry}
\usepackage[colorlinks=true,linkcolor=blue,urlcolor=blue]{hyperref}
\setlength{\emergencystretch}{2em}
\newtheorem{theorem}{Theorem}[section]
\newtheorem{lemma}[theorem]{Lemma}
\newtheorem{proposition}[theorem]{Proposition}
\newcommand{\Q}{\mathbb Q}\newcommand{\Z}{\mathbb Z}
\DeclareMathOperator{\Tr}{Tr}\DeclareMathOperator{\Nm}{Nm}
\title{Audit and extensions of the 62-coordinate quartic\\Small pole inputs, weaker hypotheses, and coefficient height}
\author{Research integration of the supplied coauthor constructions}
\date{19 November 2026}
\begin{document}\maketitle
\begin{abstract}
The supplied reduction from 68 to 62 rational coordinates passes an
independent reconstruction of its physical equations and every expanded
coefficient. We verify the local selector, its zero and tangent strata,
and the order of the integer-witness preparations. A small-input version
of the retained simultaneous-twist lemma permits deletion of a constant
from each of the two selector gates. This gives 825 monomials, or a
752-monomial weighted presentation, in the same 62 coordinates. Clearing
four reference constants by the existing unit equation gives another
presentation with 840 monomials and exact height $B+1$ on an explicitly
specified cofinal fourth-power subfamily. The supplied 63- and 64-coordinate
weak-$abc$ quartics already work at the weaker exponents $9/8$ and $17/14$.
An eight-contact extension gives a quartic for every fixed weak-$abc$
exponent below $3/2$; at exponent $7/5$ it uses 67 coordinates. All defining
claims retain the upstream elliptic, trace, prime-pattern and
simultaneous-twist inputs. The sufficient defining coefficient remains
ineffective. No reduction below 62 coordinates is proved here.
\end{abstract}

\section{Scope, arithmetic inputs, and the audited result}
All counts include the distinguished input $X$. The base is the supplied
\texttt{H10-quartics/abc62} construction, pinned to Brox repository commit
\texttt{d3a8c223716a0fc0dbf5af4eb7998d640c21048d}.
The arithmetic inputs retained from v65 are the asserted description of
$E(K)$, the simultaneous prime-pattern input with arbitrary finite
avoidance and real and local preparation, the all-point trace-index
bound, and the simultaneous congruent-twist statement stated in
Section~\ref{sec:small}. Ordinary rational $abc$ is used for the original
62-coordinate row. Each later weak-$abc$ row instead uses its single
specified exponent and the same remaining inputs. These premises are
not silently deduced from $abc$.

Write
\[
 K=\Q(\rho),\quad \rho^2=2,\quad \Delta=11-6\rho,\quad
 E:y^2=x^3-\Delta x,\quad R_0=(1+\rho,2\rho).
\]
The v65 chart theorem identifies all its nonzero-input solutions with
finite pairs $U,V\in E(K)$ of nonzero abscissa, $U\ne V$, and an allowed
reference $H$. It includes tangencies, zero ordinates, and every finite
two-torsion difference $U-V$, including $(0,0)$.
With $f(Q)=1/x(Q)$ and $S_Q=y(Q)/x(Q)$, its outputs are
\[
\begin{gathered}
s=XS_U,\qquad t=XS_V,\qquad S=S_H,\qquad T=S_{H-R_0},\\
J=X^2f(H)f(U),\qquad K'=X^2f(H-R_0)f(V).
\end{gathered}
\]
\[
 T_{\rm sc}=X^2\Tr x(U-V),\qquad
 t_1-t_0=B+X^8+T_{\rm sc}^2.
\]
Only one rational half-trace is stored. The original reference variables
are replaced invertibly by $S,T,v$, with
\[
 v(v+T-S)+\delta=0,\qquad ST+k(T-S)+2kv+2\delta=0,
\]
where $a=1+\rho$, $k=4-2\rho$, $\delta=-23+17\rho$.
The old variables are $u=(v-k-S)/a$ and $w=(v-k+T)/a$.
Their source unit identities force $u,w,v$ to be nonzero. No new
reference denominator is left unguarded.

For the two contact arguments the original weights are
\[
 w_j=\frac{\prod_{\theta\in\Theta}(t_j-\theta)}{(t_j-1)^6},\qquad
 \Theta=\{-5,-2,-1,0,2,3,4,7\}.
\]
The three contact equations are
\[
 R=t^2-2t,\qquad(R+1)Z=R^2-35R,\qquad(R+1)w=Z^2+24Z.
\]
If $t=1$, the second equation gives $0=36$; hence elimination of
$R,Z$ is valid at every rational solution. The reciprocal equations
$q_jw_j=1$ exclude every zero of the numerator. Store $q_0$ and $c$,
put $q_1=c-q_0$, and retain the literal unit equation for $c^{-1}$.

Put $C=s-S$, $I=t-T$ and
\begin{equation}\label{eq:W}
 W=1+cC^4+c^2J^2-c^3\overline I^{\,4}-c^4\overline {K'}^{\,2},
 \qquad D\Tr W=\Nm(W-1).
\end{equation}
For the supplied construction use
\begin{equation}\label{eq:gate-eps}
 2\tau_j=\varepsilon+(c-q_j)^2+q_jD^2,
 \qquad q_jb_j=\tau_j^2,\qquad \varepsilon=1.
\end{equation}
Each $b_j$ enters the native pole equations
\begin{equation}\label{eq:native}
 e^2+bp^2=1,\qquad uv=e(1-b),\qquad u^2-v^2=bq^2.
\end{equation}
The two square registers for $p^2,q^2$ are retained. The symbols $q$ in
\eqref{eq:native} and $q_j=1/w_j$ have different roles.

The replacements
\[
 A_\pm=C^2\pm c\overline I^{\,2},\qquad
 B_\pm=J\pm c\overline {K'}
\]
are invertible because $2c\ne0$, and give
$W=1+cA_+A_-+c^2B_+B_-$. Together with the supplied integral Laurent
weights and row divisors they yield precisely 50 integral quadrics in
62 rational coordinates. The weighted basis changes preserve simultaneous
vanishing. The standard sum of squares has 830 monomials and exact
height $B^2+4091$ for odd $B\ge81$; the independent check also verifies
all 833 parameter-separated terms and content one.

\section{The selector: soundness and all zero cases}
Normalize valuations by $v_p(p)=1$ and put $v(0)=+\infty$.
Outside the fixed exceptional set, every finite reciprocal-abscissa
valuation is even. Consequently $v(J)$ and $v(K')$ are even, including
when $X$ has negative valuation.

\begin{lemma}\label{lem:gate}
Let $v(q_j)=-\mu$ with $\mu>0$ odd. Suppose \eqref{eq:gate-eps} and
\eqref{eq:native} hold, where $\varepsilon=0$ or $1$. Then
\[
 v(D)>\mu/4,\qquad v(c-q_j)>-\mu/4,\qquad v(c)=-\mu.
\]
\end{lemma}
\begin{proof}
The native obstruction excludes a negative odd valuation of $b_j$.
If $\tau_j\ne0$, then
$v(b_j)=\mu+2v(\tau_j)$ is odd, so $v(\tau_j)>-\mu/2$.
For $\tau_j=0$ the same inequality holds with the adopted convention.
Any negative valuation of $\varepsilon+(c-q_j)^2$ is even; the finite
valuation of $q_jD^2$ is odd. A negative minimum cannot cancel.
Thus every negative summand valuation is greater than $-\mu/2$.
It follows that $2v(D)-\mu>-\mu/2$ and
$2v(c-q_j)>-\mu/2$. Since $-\mu<-\mu/4$, adding $q_j$ gives
$v(c)=-\mu$. Zero summands satisfy the inequalities automatically.
\end{proof}

\begin{lemma}\label{lem:select}
At every rational solution of \eqref{eq:W}, $\Tr W\ne0$.
Under the conclusions of Lemma~\ref{lem:gate}, at one place above $p$
all four raw tests $C,J,\overline I,\overline {K'}$ have positive valuation.
\end{lemma}
\begin{proof}
If $\Tr W=0$, then $\Nm(W-1)=0$. Since $K$ is a field this gives
$W=1$, contradicting $\Tr 1=2$.
If both slot values of $W$ are negative, say $a,b<0$, then
$v(\Nm(W-1))=a+b$ and $v(\Tr W)\ge\min(a,b)$.
Hence $v(D)\le\max(a,b)<0$, a contradiction. The same calculation
includes an inert place with repeated valuations. Therefore one slot
has $v(W)\ge0$. At that slot the four nonconstant summands in
\eqref{eq:W} have four different residue classes modulo four:
$-\mu,-2\mu,-3\mu,-4\mu$. This uses even valuation of $J,K'$.
A negative minimum would therefore be unique and could not cancel.
All four summands have nonnegative valuation. Since their powers of $c$
are positive, all four raw tests have strictly positive valuation.
A zero test causes no exception, and the fourth term may tie the
constant at valuation zero without affecting the proof.
\end{proof}

The v65 inverse-parameter extraction lemma now applies exactly as
stated: at every consumed contact prime, $X$ is integral and
$X(m-X)\equiv0\pmod p$, where $m$ is the index of $U-V$.
This includes vanishing individual ordinates and finite torsion
 differences. If $X=0$, soundness already holds. For the original
contact map, the unchanged ordinary-$abc$ argument has strict
absorption $602/679<1$. Thus its soundness is preserved.

\section{Small native pole inputs}\label{sec:small}
The retained simultaneous-twist statement is the following implication:
if $\beta_0>1$, $\theta\ge1$ are coprime positive squarefree integers,
both $1\pmod8$, and each is a square modulo every prime of the other,
then some positive squarefree $r$ makes $E_r$ and $E_{r\theta}$ have
infinite-order rational points whose abscissas belong to
$\beta_0\Q^{\times2}$, where $E_a:y^2=x^3-a^2x$.
This is precisely the upstream premise; it is not strengthened here.

\begin{lemma}\label{lem:small}
Under that premise, if
\[
 0<b<1,\qquad b\text{ has square reduced denominator},\qquad
 b\in\Q_2^{\times2},\qquad v_2(b)\ge4,
\]
then \eqref{eq:native} has rational witnesses with $upq\ne0$.
No congruence restriction on the odd primes in the numerator of $b$
is required.
\end{lemma}
\begin{proof}
Write $b=A/B_0^2$ in lowest terms and write
$b=\beta_0z^2$, $1-b=\theta d^2$, with $\beta_0,\theta$ positive
squarefree and $z,d\ne0$. Coprimality follows from
$\gcd(A,B_0^2-A)=1$. The two-adic assumptions give $B_0$ odd,
$v_2(A)\ge4$ even, and the odd unit part of $A$ congruent to one
modulo eight. Hence $\beta_0\equiv\theta\equiv1\pmod8$.
At a prime of $\beta_0$, $1-b\equiv1$; at a prime of $\theta$,
$b\equiv1$. Removing square factors shows that the two squarefree
integers are mutual quadratic residues as required.

If $\beta_0=1$, use $e=0$, $p=1/z$, $u=z$, $v=0$, $q=1$.
Otherwise apply the retained statement and scale the second point by
$(x,y)\mapsto(d^2x,d^3y)$. The resulting points are
$(ba_1^2,y_1)\in E_r(\Q)$ and
$(ba_2^2,y_2)\in E_{r(1-b)}(\Q)$, with all displayed nonconstant
factors nonzero. Set
\[
 e=\frac r{ba_1^2},\quad p=\frac{y_1}{b^2a_1^3},\quad
 u=\frac{a_2}{a_1},\quad v=\frac{e(1-b)}u,\quad
 q=\frac{y_2}{b^2a_1^3u^2}.
\]
The two curve equations give respectively $e^2+bp^2=1$ and
$u^2-v^2=bq^2$; the middle equation is the definition of $v$.
All divisions are by nonzero numbers, and the infinite-order points
have nonzero ordinates. These identities are also checked symbolically
in the independent verifier.
\end{proof}

\section{Constant-free gates and completeness}
We now take $\varepsilon=0$ in \eqref{eq:gate-eps}. Lemmas
\ref{lem:gate} and~\ref{lem:select} prove the unchanged soundness.
Fix an integer $X\ne0$. As in v65, first fix $A=[X]R_0$,
its scaled trace, the positive gap, and bounded real and fixed-local
neighborhoods of the reference and numerator point functions.
All contact lists and their avoidance values are fixed before prime
selection. The two translated lists are disjoint because the gap is
at least 81 and exceeds the contact diameter.

At a selected simple contact for block $j$, $v(w_j)=1$ and the other
reciprocal $q_{1-j}$ is a unit after excluding the numerator and
denominator primes of its fixed nonzero contact value. The successful
slot has $v(C),v(\overline I)\ge3$ and $v(J),v(\overline {K'})\ge4$.
Thus $W\equiv1$ there. At the other slot the reciprocal products are
units; the $-c^4\overline {K'}^{\,2}$ term has valuation $-4$ and
is the unique minimum. Thus $v(D)>0$. The new gate gives
$2\tau_j\equiv q_{1-j}^2\not\equiv0\pmod p$, and hence
$v(b_j)=1$. For the supplied $\varepsilon=1$ gates one additionally
avoids the primes of $1+q_{1-j}^2$; their values are positive rational
numbers, so this is legitimate finite avoidance. Both preparations
therefore respect the required order of choices.

As the two real arguments tend to positive infinity with fixed gap,
\[
 w_j\sim t_j^2,\quad q_j\sim t_j^{-2},\quad c=O(t_j^{-2}),\quad
 W-1=O(t_j^{-2}),\quad D=O(t_j^{-4}),
\]
\[
 \tau_j\sim\frac1{2t_j^4},\qquad b_j\sim\frac1{4t_j^6}.
\]
These bounds are uniform on the already fixed compact point
neighborhoods. Thus $0<b_j<1$. The same asymptotics at fixed places
are obtained by increasing the prescribed slope-denominator depth.
At two, if $v_2(q_j)=2k$ at the common leading order, then eventually
$v_2(\tau_j)=4k-1$ and $v_2(b_j)=6k-2\ge4$.
The contact squareclass formula makes $w_j$, and consequently $b_j$,
a two-adic square. Its selected primes have positive odd valuations;
every other valuation is even because the gate contributes a square.
Thus the reduced denominator is a square. Lemma~\ref{lem:small}
supplies both actual pole blocks.

For $X=0$ use the v65 zero fiber: $H=2R_0$, $s=S$, $t=T$,
$F=G=d=J=K'=0$, and $h=t^2/(s-t)$. Its denominator is nonzero
because $\Nm(S-T)=7199/328$. The half-trace gives a fixed positive gap.
Now every raw test is zero, so $W=1$, $D=0$ and
$\tau_j=q_{1-j}^2/2$. The same prime, local, and real preparation gives
both small native inputs. This proves completeness without deleting
any zero-input, tangent, or zero-ordinate branch.

\section{Coefficients and exact height}
In stored coordinates the only change from the supplied system is
removing $-x_{35}x_{62}$ from \texttt{gate0} and
$-x_{49}x_{62}$ from \texttt{gate1}. The same 62 coordinates, 50 rows,
Laurent weights, and inverse maps remain valid. Direct expansion gives
825 monomials and height $B^2+4091$. On the supplied cofinal scaling
$B=n^4$, the balanced table has 825 monomials and height $B+4091$
for odd $n\ge3$. Its supplied positive weighted row operations give
752 monomials and height $B+758541$ in the same range.

For the smaller height presentation start from the balanced rows.
Its unit row is $r_{50}=n x_{61}x_{62}-1$. The constants of the four
reference rows are respectively $23,17,46,34$. Replace those rows by
\[
 r_1+23r_{50},\quad r_2+17r_{50},\quad
 r_3+46r_{50},\quad r_4+34r_{50}.
\]
These are invertible row operations. The resulting sum of 50 squares
has 840 coordinate monomials and 841 parameter-separated terms.
The maximum absolute coefficients at powers $n^0,n^1,n^2,n^3,n^4$
are respectively
\[
                    1548,\quad1280,\quad4091,\quad4,\quad1.
\]
The sole mixed profile is its constant coefficient $n^4+1$.
Every other coordinate monomial has a single parameter power.
Consequently, for every positive integer $n$,
\[
                 H=\max(4091n^2,n^4+1).
\]
For odd $n\ge65$ this is exactly $B+1$.
Every coordinate occurs, the degree is exactly four, and the coefficient
of $X^4$ is one, proving content one. The cutoff $65$ certifies the
algebraic height formula only; it is not a sufficient defining shift.

\section{Fixed weak-abc hypotheses and new quartic families}
Assume weak rational $abc$ at one fixed exponent $\kappa>1$.
For six contacts use
\[
 \Theta_6=\{1,31,37,41,46,61\}.
\]
The map $z=2+60/(t-61)$ sends them to
$\{0,-2,-1/2,1,-1,\infty\}$. The identity
\[
 z(z+2)^3-(2z+1)^3=(z-1)^3(z+1)
\]
gives the six-contact rate $\gamma_6=8/\kappa-6$.
For eight contacts use
\[
 \Theta_8=\{1,8,22,28,29,36,43,64\}.
\]
The identity
\[
 (v+4)(v+20)^3(v-5)^2-(v-4)(v-20)^3(v+5)^2=108v^5
\]
and the map $t=43+210/(v-10)$ give the eight-contact rate
$\gamma_8=12/\kappa-8$.
Indeed the coprime degree-four or degree-six homogeneous abc triple
has bounded common divisor, so its radical has logarithm at least
$(k-2)h(t)/\kappa-O(1)$. Its full homogeneous contact form has degree
$k$. The inequality $S_1(N)\ge\operatorname{rad}(N)^2/|N|$, where
$S_1$ is the product of primes occurring exactly once, gives
\begin{equation}\label{eq:rate}
 \log M(t)\ge\gamma_kh(t)-O(1),\qquad
 \gamma_k=\frac{2(k-2)}\kappa-k.
\end{equation}
Here $M$ can include all positive odd contact primes. The omitted
primes form a fixed finite set consisting of the point-chart bad primes,
contact coefficients, differences and transformation resultants.
All constants are independent of the later coefficient $B$.

Put $G_k(t)=\prod_{\theta\in\Theta_k}(t-\theta)$ and
$w(t)=G_k(t)/t^{k-2}$. Since zero is not a contact, the equations
$te=1$, $f=e^2$ and, when $k=8$, $g=ef$ give quadratic charts:
\[
\begin{split}
 k=6:\quad w={}&t^2+a_5t+a_4+a_3e+a_2f+a_1ef+a_0f^2,\\
 k=8:\quad w={}&t^2+a_7t+a_6+a_5e+a_4f+a_3g+a_2f^2+a_1fg+a_0g^2,
\end{split}
\]
where $G_k(t)=\sum a_it^i$. Each six-contact block uses three
helpers, as before; each eight-contact block uses four.
Under the existing Laurent weights all these coordinates have physical
factor $c$, and division by $c^2$ makes the displayed equations literal
quadrics using $c^{-1}$. At $t=0$ the inverse equation is impossible;
at every contact $w=0$ contradicts $wq=1$.

Replace the gap by
\[
             G=t_1-t_0=B+X^{2^r}+T_{\rm sc}^2,\qquad r\ge3.
\]
Keep the original $X^2,X^4$ registers and add successive squares up to
$X^{2^{r-1}}$. Thus the respective coordinate counts are
\begin{equation}\label{eq:count}
                  N_6(r)=59+r,\qquad N_8(r)=61+r.
\end{equation}
There are $N_k(r)-12$ rational quadratic rows.

\begin{theorem}\label{thm:weak}
If $\gamma_k>0$ and
\begin{equation}\label{eq:strict}
 \bigl(2^r-\kappa(2^{r-1}+2)\bigr)\gamma_k>4(2\kappa-1),
\end{equation}
then, under the retained remaining premises, every sufficiently large
fixed odd $B$ gives integer input projection for the above quartic.
In particular $k=6$ covers every fixed $1<\kappa<4/3$ for sufficiently
large $r$, and $k=8$ covers every fixed $1<\kappa<3/2$.
\end{theorem}
\begin{proof}
The local selector is unchanged. For $X=u/v\ne0$ in lowest terms and
$X\ne m$, both contact prime products divide
$|u(mv-u)|\le(1+|m|)\exp(2h(X))$. Equation~\eqref{eq:rate} gives
\[
 h(G)\le\frac2{\gamma_k}\bigl(2h(X)+\log(1+|m|)\bigr)+O(1).
\]
Set $s=X^{2^r}+T_{\rm sc}^2\ge0$. Then $h(s)\le h(G)$ and
$h(G)\ge\log B$. The retained signed mixed-power height theorem,
applied with the plus sign, gives
\[
 h(X)\le\frac{2\kappa-1}{2^r-\kappa(2^{r-1}+2)}h(s)+O(1).
\]
Condition~\eqref{eq:strict} makes the resulting coefficient of $h(G)$
strictly smaller than one. Thus $h(G)$ and $h(X)$ are bounded by fixed
multiples of $\log(1+|m|)+1$. Also
$2h(T_{\rm sc})\le h(s)+2^rh(X)+\log2$.
The retained trace-index bound
$m^2\le C(h(T_{\rm sc})+2h(X))+J_0$ bounds $m$, including finite
torsion differences by taking $m=0$. It follows that $h(G)$ is uniformly
bounded independently of $B$, contradicting $h(G)\ge\log B$.

For completeness, the point family and the constant-free gates are
prepared as above. Both fixed contact lists are disjoint because their
diameter is at most 63 and the gap is at least 81. For
$t_j=N/L+g_j$ the prime-pattern theorem provides the positive prime values
$p_{j,\theta}=N+(g_j-\theta)L$, with all specified avoidance and local
conditions. The exact squareclass formula is
\[
 w_j=\frac{\prod_{\theta\in\Theta_k}p_{j,\theta}}
             {L^2(N+g_jL)^{k-2}}.
\]
Its denominator is a square and its leading real coefficient is one.
The same selected-prime, real, and two-adic calculations therefore give
the small native pole inputs, including the separate zero-input fiber.
The new power registers are determined uniquely. This proves both
implications.
\end{proof}

For $k=6,r=4$, the exact strict coefficient at $\kappa=9/8$ is
$18/19$; for $k=6,r=5$ at $\kappa=17/14$ it is $68/71$.
These conclusions also strengthen the hypotheses of the supplied
63- and 64-coordinate fractional-contact families without changing any
of their coefficients. Their originally stated exponents were $16/15$
and $121/100$. For $k=8,r=6$, exponent $7/5$ has coefficient $63/82$.

The new integer-contact, constant-free families have the following
independently verified presentations, all with $B=n^4$:
\begin{center}
\begin{tabular}{rrrrrr}
\toprule
$\kappa$&Variables&Rows&Monomials&Height&Odd algebraic cutoff\\
\midrule
$9/8$&63&51&863&$B+1$&$n\ge17207$\\
$17/14$&64&52&866&$B+1$&$n\ge17207$\\
$7/5$&67&55&905&$B+1$&$n\ge192405$\\
\bottomrule
\end{tabular}
\end{center}
The first two coefficient caps at powers $n^0$ through $n^4$ are
\[
 39067747815195764,\quad64245567274,\quad296062522,\quad434,\quad1;
\]
for the third they are
\[
 524140711592652177408,\quad8551459564032,\quad37019305472,\quad462,\quad1.
\]
The sole mixed coefficient is $n^4+1$, so the stated height cutoffs
follow directly and hold uniformly thereafter. These are algebraic
height cutoffs, not effective integer-defining thresholds.

\section{A number-field version of the shifted-norm selector}
\begin{proposition}
Let $L/\Q$ be a number field of degree $d$, and let
$e_{d-1}(W)$ be the penultimate elementary symmetric polynomial in its
$d$ conjugates. The equation
\[
                D e_{d-1}(W)=\Nm_{L/\Q}(W-1)
\]
forces $e_{d-1}(W)\ne0$ at every rational solution with $W\in L$.
If all local conjugate valuations of $W$ are negative, then
$v(D)\le\max_\sigma v(\sigma W)<0$. Conversely, suppose $k\ge1$
conjugates satisfy $W_\sigma\equiv1$ and the others are poles.
If $p\nmid k$, then $v(D)>0$, with $D=0$ allowed.
\end{proposition}
\begin{proof}
A zero denominator would force $\Nm(W-1)=0$, hence $W=1$, whereas
$e_{d-1}(1)=d\ne0$. Put $a_\sigma=v(\sigma W)$. If all are negative,
the numerator has valuation $\sum a_\sigma$ and the denominator has
valuation at least $\sum a_\sigma-\max a_\sigma$, proving the bound.
For the converse divide numerator and denominator by the product of
the pole conjugates. The denominator reduces to $k$ modulo the maximal
ideal: precisely the terms omitting one of the successful conjugates
survive. The numerator contains the product of their positive-order
errors $W_\sigma-1$ and unit factors from all pole conjugates. Thus its
valuation is positive, or it is zero identically. Excluding the finitely
many primes at most $d$ makes the condition automatic for every $k$.
\end{proof}
The residue-class separation argument similarly extends to $2^a$ tests
when their powered valuations are multiples of $2^a$ and the selector
unit has negative odd valuation. Distinct exponents modulo $2^a$ give
a unique negative minimum. This field theorem concerns the arithmetic
selector; in higher degree its norm equation need not remain quadratic.
It is not a new quartic coordinate-count claim.

\section{Audit record and limitations}
The supplied standard-library checker and its independent SymPy checker
both pass. The latter compares the full five-row span containing the
reference equations and a pole row with the pinned v65 source, reverses
every Laurent and hyperbolic substitution, and independently expands
all 50 squared quadrics. The supplementary fractional-contact 63- and
64-coordinate checkers also pass. The new independent checks verify the
changed gate identity, the two native twist witness identities, both
Belyi identities and transformations, every new contact and power row,
all scaling and row operations, every expanded coefficient, support,
content, and exact coefficient caps.

The computation is not a proof of the external arithmetic premises or
an effective defining coefficient. Neither the recovered charts nor
finite valuation searches establish a smaller rational coordinate count.
Several direct eliminations merely replace a field helper by another:
for example removing the secant helper $d$ changes the scale equation
into $\Delta G^2+t^2G=X^4$, which needs a square helper to stay quadratic.
No optimality theorem is inferred from these unsuccessful reductions.

\begin{thebibliography}{9}\raggedright
\bibitem{v65} J.~Brox and contributors, \emph{A quartic in 68 variables from
translated secants}, 17 November 2026, repository file
\nolinkurl{papers/quartic_existence/quartic_abc68_20261117_v65.tex}.
\bibitem{v58} J.~Brox and contributors, signed mixed-power height theorem,
repository file
\nolinkurl{papers/elliptic_heights/signed_mixed_power_height_20261117_v58.tex}.
\bibitem{v22} J.~Brox and contributors, direct-input weak-$abc$ construction,
repository file
\nolinkurl{papers/conditional_few_variables/direct_input13_20261110_v22.tex}.
\bibitem{coauthor} Supplied coauthor bundles
\texttt{H10-quartic-constructions.zip} and
\texttt{H10-eleven-variable-constructions.zip}, particularly
\texttt{abc62/PROOF.md}, \texttt{counts/abc/quartics/proof.md}, and
\texttt{eight\_cusps.py}. The independent checks certify relative
transfer from their pinned upstream, with the arithmetic scope stated above.
\end{thebibliography}
\end{document}
